% Автоматаар үүсгэсэн — эх файл: edge/README.md
\chapter{Ирмэгийн давхарга — лавлах}

\begin{quote}
СТОП: Энэ фолдерыг \textbf{зөөврийн компьютер дээр ажиллуулахгүй}. Компьютерийн хэсэг нь \texttt{../stack/}.
\end{quote}

Pi 3B: \textbf{1 GB RAM} (\textasciitilde925 MiB харагдана, \textasciitilde720--780 MiB ашиглах боломжтой), 4×Cortex-A53 @1.2 GHz, 100 Mbit (USB 2.0-ийн зурвасыг хуваана), microSD.

\hypertarget{ux445ux443ux440ux434ux430ux43d-ux44dux445ux43bux44dux445}{%
\section{Хурдан эхлэх}\label{ux445ux443ux440ux434ux430ux43d-ux44dux445ux43bux44dux445}}

\begin{Shaded}
\begin{Highlighting}[]
\FunctionTok{cp}\NormalTok{ .env.example .env}
\FunctionTok{nano}\NormalTok{ .env            }\CommentTok{\# CLOUD\_HOST = зөөврийн компьютерийн LAN IP (localhost БИШ)}
\FunctionTok{make}\NormalTok{ bridge          }\CommentTok{\# .env → mosquitto/conf.d/bridge.conf}
\FunctionTok{make}\NormalTok{ up              }\CommentTok{\# mosquitto асаана}
\FunctionTok{make}\NormalTok{ link            }\CommentTok{\# → bridge/state 1 гарвал холбогдсон}
\FunctionTok{make}\NormalTok{ agent           }\CommentTok{\# ирмэгийн агент (өөр терминалд)}
\end{Highlighting}
\end{Shaded}

\hypertarget{ux44eux443-ux44dux43dux434-ux430ux436ux438ux43bux43bux430ux445-ux432ux44d}{%
\section{Юу энд ажиллах вэ}\label{ux44eux443-ux44dux43dux434-ux430ux436ux438ux43bux43bux430ux445-ux432ux44d}}

\begingroup\scriptsize\begin{longtable}[]{@{}
  >{\raggedright\arraybackslash}p{(\columnwidth - 6\tabcolsep) * \real{0.2500}}
  >{\raggedright\arraybackslash}p{(\columnwidth - 6\tabcolsep) * \real{0.2500}}
  >{\raggedright\arraybackslash}p{(\columnwidth - 6\tabcolsep) * \real{0.2500}}
  >{\raggedright\arraybackslash}p{(\columnwidth - 6\tabcolsep) * \real{0.2500}}@{}}
\toprule
\begin{minipage}[b]{\linewidth}\raggedright
Хэсэг
\end{minipage} & \begin{minipage}[b]{\linewidth}\raggedright
Хэрхэн
\end{minipage} & \begin{minipage}[b]{\linewidth}\raggedright
RAM
\end{minipage} & \begin{minipage}[b]{\linewidth}\raggedright
Үүрэг
\end{minipage} \\
\midrule
\endhead
\bottomrule
\endlastfoot
\texttt{mosquitto} & Docker & 12--25 MiB & локал брокер + үүл рүү гүүр + диск дээрх дараалал \\
ирмэгийн агент & venv (Docker биш) & 45--70 MiB & Pi-гийн бодит хэмжүүр + процессын дохио → UNS \\
TFLite дүгнэлт & агентын дотор & +40--90 MiB & Лаб 6: локал аномали илрүүлэлт \\
K3s & Лаб 8 & 250--350 MiB & Docker-ийг \textbf{зогсоож} ажиллуулна \\
\end{longtable}\endgroup

Агентыг Docker-т биш \textbf{venv-д} ажиллуулж байгаа шалтгаан: контейнерийн давхарга Pi 3B дээр 60--80 MiB нэмнэ, мөн кодыг засаад шууд дахин ажиллуулах нь илүү хурдан. Лаб 8-д K3s-д оруулахын тулд \texttt{agent/Dockerfile}-ээр дүрс барина.

\hypertarget{make-ux43aux43eux43cux430ux43dux434ux443ux443ux434}{%
\section{\texorpdfstring{\texttt{make} командууд}{make командууд}}\label{make-ux43aux43eux43cux430ux43dux434ux443ux443ux434}}

\begingroup\scriptsize\begin{longtable}[]{@{}
  >{\raggedright\arraybackslash}p{(\columnwidth - 2\tabcolsep) * \real{0.5000}}
  >{\raggedright\arraybackslash}p{(\columnwidth - 2\tabcolsep) * \real{0.5000}}@{}}
\toprule
\begin{minipage}[b]{\linewidth}\raggedright
Комманд
\end{minipage} & \begin{minipage}[b]{\linewidth}\raggedright
Үйлдэл
\end{minipage} \\
\midrule
\endhead
\bottomrule
\endlastfoot
\texttt{make\ bridge} & \texttt{.env}-ээс \texttt{conf.d/bridge.conf} үүсгэнэ (\textbf{эхлээд энэ}) \\
\texttt{make\ up} & \texttt{bridge} + mosquitto асаана \\
\texttt{make\ link} & гүүрний төлөв сонсоно (1 = холбогдсон) \\
\texttt{make\ agent} & ирмэгийн агент (venv байхгүй бол өөрөө үүсгэнэ) \\
\texttt{make\ mem} & санах ой, throttle, температур, ачаалал \\
\texttt{make\ check} & бэлэн байдлын шалгалт (swap, cgroup, NTP, bridge.conf) \\
\texttt{make\ logs} & mosquitto-гийн лог \\
\texttt{make\ ps} & контейнер + \texttt{docker\ stats} \\
\texttt{make\ down} & зогсоох (өгөгдөл үлдэнэ) \\
\texttt{make\ clean} & зогсоох + volume + \texttt{bridge.conf} устгах \\
\end{longtable}\endgroup

\hypertarget{ux444ux430ux439ux43bux443ux443ux434}{%
\section{Файлууд}\label{ux444ux430ux439ux43bux443ux443ux434}}

\begin{verbatim}
edge/
├── docker-compose.yml            mosquitto (+ containerized профайл, Лаб 8)
├── .env / .env.example           CLOUD_HOST, UNS байрлал, агентын тохиргоо
├── Makefile
├── mosquitto/
│   ├── mosquitto.conf            persistence, дараалал, хязгаарлалт
│   └── conf.d/
│       ├── bridge.conf.template  ЗАГВАР (Git-д байна)
│       └── bridge.conf           ҮҮСГЭСЭН (.gitignore-д — IP агуулна)
└── agent/
    ├── edge_agent.py             ирмэгийн агент
    ├── requirements.txt
    ├── Dockerfile                Лаб 8-д K3s-д хэрэгтэй
    └── cnc302-edge-agent.service systemd (тэжээл тасрахад сэргэнэ)
\end{verbatim}

\hypertarget{edge_agent.py-ux433ux43eux43b-ux441ux43eux43dux433ux43eux43bux442ux443ux443ux434}{%
\section{\texorpdfstring{\texttt{edge\_agent.py} --- гол сонголтууд}{edge\_agent.py --- гол сонголтууд}}\label{edge_agent.py-ux433ux43eux43b-ux441ux43eux43dux433ux43eux43bux442ux443ux443ux434}}

\begin{Shaded}
\begin{Highlighting}[]
\ExtensionTok{.venv/bin/python}\NormalTok{ agent/edge\_agent.py }\AttributeTok{{-}{-}help}
\end{Highlighting}
\end{Shaded}

\begingroup\scriptsize\begin{longtable}[]{@{}
  >{\raggedright\arraybackslash}p{(\columnwidth - 2\tabcolsep) * \real{0.5000}}
  >{\raggedright\arraybackslash}p{(\columnwidth - 2\tabcolsep) * \real{0.5000}}@{}}
\toprule
\begin{minipage}[b]{\linewidth}\raggedright
Тохиргоо
\end{minipage} & \begin{minipage}[b]{\linewidth}\raggedright
Утга
\end{minipage} \\
\midrule
\endhead
\bottomrule
\endlastfoot
\texttt{-\/-dry-run} & брокергүйгээр зөвхөн хэвлэнэ (SETUP-д шалгахад) \\
\texttt{-\/-interval\ 2.0} & телеметрийн үе (сек) \\
\texttt{-\/-qos\ 0\textbackslash{}\textbar{}1\textbackslash{}\textbar{}2} & нийтлэх QoS \\
\texttt{-\/-anomaly-rate\ 0.02} & санамсаргүй аномали оруулах магадлал \\
\texttt{-\/-detector\ \textless{}model\textgreater{}} & TFLite загвар (Лаб 6). Байхгүй бол босгын горим \\
\texttt{-\/-threads\ 2} & дүгнэлтийн урсгал. 4×A53-д 4 нь ихэвчлэн 2-оос \textbf{удаан} \\
\texttt{-\/-filter} & зөвхөн аномали + үе үе хураангуй илгээнэ (уплинк хэмнэнэ) \\
\texttt{-\/-health-every\ 15} & хэдэн мөчлөг тутам эрүүл мэндийн мессеж \\
\end{longtable}\endgroup

Дохио нь \texttt{DEVICE\_ID}-аар seed хийгдсэн тул \textbf{давтагдана} --- нэг Pi үргэлж ижил цуваа өгнө. Лаб 6-д загвар харьцуулахад зайлшгүй.

\hypertarget{ux431ux430ux439ux43dux433ux430-ux430ux436ux438ux43bux43bux443ux443ux43bux430ux445-ux441ux43eux43dux433ux43eux43bux442}{%
\section{Байнга ажиллуулах (сонголт)}\label{ux431ux430ux439ux43dux433ux430-ux430ux436ux438ux43bux43bux443ux443ux43bux430ux445-ux441ux43eux43dux433ux43eux43bux442}}

\begin{Shaded}
\begin{Highlighting}[]
\FunctionTok{sudo}\NormalTok{ cp agent/cnc302{-}edge{-}agent.service /etc/systemd/system/}
\FunctionTok{sudo}\NormalTok{ systemctl daemon{-}reload}
\FunctionTok{sudo}\NormalTok{ systemctl enable }\AttributeTok{{-}{-}now}\NormalTok{ cnc302{-}edge{-}agent}
\ExtensionTok{journalctl} \AttributeTok{{-}u}\NormalTok{ cnc302{-}edge{-}agent }\AttributeTok{{-}f}
\end{Highlighting}
\end{Shaded}

\texttt{MemoryMax=200M} тавьсан --- агент хэтэрвэл systemd таслана. Pi-г нэг GB-тай гэдгийг мартаж болохгүй.

\hypertarget{ux445ux44dux43cux436ux438ux43bux442ux438ux439ux43d-ux4e9ux43cux43dux4e9-ux437ux430ux430ux432ux430ux43b}{%
\section{Хэмжилтийн өмнө ЗААВАЛ}\label{ux445ux44dux43cux436ux438ux43bux442ux438ux439ux43d-ux4e9ux43cux43dux4e9-ux437ux430ux430ux432ux430ux43b}}

\begin{Shaded}
\begin{Highlighting}[]
\ExtensionTok{pkill} \AttributeTok{{-}f}\NormalTok{ vscode{-}server        }\CommentTok{\# VS Code Remote сервер 150–250 MiB иднэ}
\ExtensionTok{vcgencmd}\NormalTok{ get\_throttled        }\CommentTok{\# 0x0 байх ЁСТОЙ}
\FunctionTok{free} \AttributeTok{{-}m} \KeywordTok{|} \FunctionTok{awk} \StringTok{\textquotesingle{}/Mem:/\{print $7 " MiB available"\}\textquotesingle{}}
\end{Highlighting}
\end{Shaded}

\texttt{get\_throttled} нь \texttt{0x0} биш бол тэр хэмжилт \textbf{хүчингүй}. Хөргөөд дахин хий.

\hypertarget{ux442ux4afux433ux44dux44dux43cux44dux43b-ux430ux43bux434ux430ux430}{%
\section{Түгээмэл алдаа}\label{ux442ux4afux433ux44dux44dux43cux44dux43b-ux430ux43bux434ux430ux430}}

\begingroup\scriptsize\begin{longtable}[]{@{}
  >{\raggedright\arraybackslash}p{(\columnwidth - 2\tabcolsep) * \real{0.5000}}
  >{\raggedright\arraybackslash}p{(\columnwidth - 2\tabcolsep) * \real{0.5000}}@{}}
\toprule
\begin{minipage}[b]{\linewidth}\raggedright
Шинж
\end{minipage} & \begin{minipage}[b]{\linewidth}\raggedright
Шийдэл
\end{minipage} \\
\midrule
\endhead
\bottomrule
\endlastfoot
\texttt{bridge/state\ 0} & \texttt{CLOUD\_HOST} буруу, эсвэл компьютерийн галт хана (SETUP.md А.6) \\
Гүүр холбогдсон ч мессеж алга & \texttt{SITE} нь \texttt{.env} ба \texttt{bridge.conf}-д зөрсөн → \texttt{make\ bridge\ \&\allowbreak{}\&\allowbreak{}\ make\ restart} \\
\texttt{.env\ олдсонгүй} & \texttt{cp\ .env.example\ .env} \\
\texttt{paho-mqtt\ суулгаагүй} & \texttt{make\ venv} \\
Контейнер дахин дахин эхэлнэ & OOM. \texttt{docker\ inspect\ …\ OOMKilled}, \texttt{free\ -m}, swap шалга \\
\texttt{exec\ format\ error} & 32-bit OS. \texttt{uname\ -m} → \texttt{aarch64} байх ёстой \\
\end{longtable}\endgroup

Дэлгэрэнгүйг \texttt{.\allowbreak{}.\allowbreak{}/\allowbreak{}docs/\allowbreak{}troubleshooting.\allowbreak{}md}-ээс.
